% Operational definitions of the three self-improvement mechanisms (feedback, memory, search).
\begin{table*}[t]
\centering
\footnotesize
\renewcommand{\arraystretch}{1.3}
\begin{tabular}{@{}l p{0.42\textwidth} p{0.30\textwidth}@{}}
\toprule
\textbf{Mechanism} & \textbf{Counts (a system has it if and only if\ldots)} & \textbf{Does not count} \\
\midrule
\textbf{Feedback (F)} &
an execution-grounded signal is read and used to revise the code \emph{within} a task: a
failure detector firing, a sensory or textual summary of what went wrong, or a verifier or
reward obtained by running the program (a simulated reward qualifies when the agent uses it to
revise) &
a fresh human instruction; a fixed pre-trained value function that is never queried at run
time; a natural-language critique that is not grounded in execution \\
\textbf{Memory (M)} &
content written \emph{at run time} is stored across task boundaries and retrieved for later
tasks: a persistent code-skill library, or a retrievable episodic-text memory of past
solutions &
the model's frozen pre-trained weights; a within-task scratchpad discarded at episode end; a
fixed API library the agent cannot extend \\
\textbf{Search (S)} &
more than one candidate program is maintained, scored by execution, and selected or mutated
among: sampled-program selection, evolutionary mutation, or beam search &
a single sequential chain of repairs on one candidate (that is Feedback); one-shot sampling
with no selection; repeated prompting whose candidates are never compared \\
\bottomrule
\end{tabular}
\caption{Operational definitions of the three mechanisms used to place a system on the
code-as-policy ladder (\S\ref{sec:code}). These resolve the ambiguous boundary cases: a
failure detector or a grounded critique counts as Feedback, trained weights do not count as
Memory, and repeated prompting counts as Search only when its candidates are compared and
selected.}
\label{tab:defs}
\end{table*}
